% !TEX TS-program = pdflatex
% Example: a Politecnico di Milano master's thesis in the "classical"
% format, using the modular academicthesis class + academicthesis-polimi.
% Mirrors the official Scuola 3I template (Thesis_classical_format).
% Finds the packages in ../core, ../languages, ../universities (any editor).
% academicthesis-repo-paths.tex -- lets showcases find the packages in this
% repository with ANY editor or engine (Texifier/Texpad, TeXShop, TeXstudio,
% plain pdflatex), not only with latexmk.
%
% \input it on the first line of a showcase, BEFORE \documentclass:
%
%   % academicthesis-repo-paths.tex -- lets showcases find the packages in this
% repository with ANY editor or engine (Texifier/Texpad, TeXShop, TeXstudio,
% plain pdflatex), not only with latexmk.
%
% \input it on the first line of a showcase, BEFORE \documentclass:
%
%   % academicthesis-repo-paths.tex -- lets showcases find the packages in this
% repository with ANY editor or engine (Texifier/Texpad, TeXShop, TeXstudio,
% plain pdflatex), not only with latexmk.
%
% \input it on the first line of a showcase, BEFORE \documentclass:
%
%   \input{academicthesis-repo-paths}
%   \documentclass[university=luiss]{academicthesis}
%
% It only extends LaTeX's own file search list (\input@path), which
% \documentclass, \RequirePackage, \input, \IfFileExists and
% \includegraphics all consult after TeX's normal search. The packages
% themselves stay path-free and are still loaded purely by name. Paths are
% relative to this showcases/ folder.
%
% \input@path does not search subfolders, so add one line per new
% university folder here (languages/ and core/ are flat).
\makeatletter
\providecommand*\input@path{}
\edef\input@path{%
  \input@path
  {../core/}%
  {../languages/}%
  {../universities/luiss/}%
  {../universities/polimi/}%
  {../universities/bocconi/}%
}
\makeatother
\endinput

%   \documentclass[university=luiss]{academicthesis}
%
% It only extends LaTeX's own file search list (\input@path), which
% \documentclass, \RequirePackage, \input, \IfFileExists and
% \includegraphics all consult after TeX's normal search. The packages
% themselves stay path-free and are still loaded purely by name. Paths are
% relative to this showcases/ folder.
%
% \input@path does not search subfolders, so add one line per new
% university folder here (languages/ and core/ are flat).
\makeatletter
\providecommand*\input@path{}
\edef\input@path{%
  \input@path
  {../core/}%
  {../languages/}%
  {../universities/luiss/}%
  {../universities/polimi/}%
  {../universities/bocconi/}%
}
\makeatother
\endinput

%   \documentclass[university=luiss]{academicthesis}
%
% It only extends LaTeX's own file search list (\input@path), which
% \documentclass, \RequirePackage, \input, \IfFileExists and
% \includegraphics all consult after TeX's normal search. The packages
% themselves stay path-free and are still loaded purely by name. Paths are
% relative to this showcases/ folder.
%
% \input@path does not search subfolders, so add one line per new
% university folder here (languages/ and core/ are flat).
\makeatletter
\providecommand*\input@path{}
\edef\input@path{%
  \input@path
  {../core/}%
  {../languages/}%
  {../universities/luiss/}%
  {../universities/polimi/}%
  {../universities/bocconi/}%
}
\makeatother
\endinput

\documentclass[university=polimi, master, format=classical, language=english]{academicthesis}

% The class only loads what the PoliMi layout needs; anything else is
% your choice. These are the ones the official template example uses.
\usepackage{amssymb}
\usepackage{bm}
\usepackage{float}    % [H] placement
\usepackage{colortbl} % \rowcolor{bluepoli!40}
\usepackage[colorlinks=true, linkcolor=black, citecolor=black, urlcolor=black]{hyperref}

\thesissetup{
  title         = {Title},
  candidate     = {Name Surname},
  idnumber      = {000000},
  degree        = {Xxxxxxx Engineering - Ingegneria Xxxxxxx},
  supervisor    = {Prof. Name Surname},
  cosupervisors = {Name Surname, Name Surname},
  academicyear  = {20XX-XX},
}

%\addbibresource{example.bib}

\begin{document}

\frontpage

\frontmatter

\begin{abstractthesis}[keywords = {here, the keywords, of your thesis}]
  Here goes the Abstract in English of your thesis followed by a list of
  keywords. The Abstract is a concise summary of the content of the thesis
  (single page of text) and a guide to the most important contributions
  included in your thesis.

  It should be a self-contained text and should be clear to someone who
  hasn't (yet) read the whole manuscript.
\end{abstractthesis}

\begin{abstractthesis}[language = italian,
                       keywords = {qui, vanno, le parole chiave, della tesi}]
  Qui va l'Abstract in lingua italiana della tesi seguito dalla lista di
  parole chiave.
\end{abstractthesis}

\tableofcontents

\mainmatter

\chapter*{Introduction}
This document is an example of the PoliMi classical format built with the
modular \texttt{academicthesis} class. The cover page contains the title,
the study programme, the author, the student ID, the advisor, the
co-advisors and the academic year, all set with \verb|\thesissetup|.

An unnumbered chapter like this one is created with \verb|\chapter*{...}|
and still appears in the table of contents and in the running header.

\chapter{Chapter one}
\label{ch:chapter_one}
In this chapter additional useful information are reported.

\section{Sections and subsections}
Chapters are typically subdivided into sections and subsections, and,
optionally, subsubsections, paragraphs and subparagraphs.

\subsection{A subsection}
Sections and subsections are numbered and coloured.

\section{Equations}
Maxwell's equations read:
\begin{align}
  \nabla\cdot \bm{D} & = \rho, \label{eq:maxwell1} \\
  \nabla \times \bm{E} + \frac{\partial \bm{B}}{\partial t} & = \bm{0}, \\
  \nabla\cdot \bm{B} & = 0, \\
  \nabla \times \bm{H} - \frac{\partial \bm{D}}{\partial t} & = \bm{J}.
\end{align}

\section{Figures and tables}
\begin{figure}[H]
  \centering
  \includegraphics[width=0.3\textwidth]{logopolimi3i}
  \caption{Caption of the Figure to appear in the List of Figures.}
  \label{fig:logo}
\end{figure}

\begin{table}[H]
  \centering
  \begin{tabular}{|p{3em} c c c|}
    \hline
    \rowcolor{bluepoli!40}
    & \textbf{column 1} & \textbf{column 2} & \textbf{column 3} \\
    \hline \hline
    \textbf{row 1} & 1 & 2 & 3 \\
    \textbf{row 2} & $\alpha$ & $\beta$ & $\gamma$ \\
    \hline
  \end{tabular}
  \caption{Caption of the Table to appear in the List of Tables.}
  \label{table:example}
\end{table}

\section{Theorems and propositions}
\begin{theorem}
  Write here your theorem.
\end{theorem}
\textit{Proof.} If useful you can report here the proof.

\begin{proposition}
  Write here your proposition.
\end{proposition}

\chapter{Conclusions and future developments}
A final chapter containing the main conclusions of your research and
possible future developments of your work.

%\printbibliography

\appendix
\chapter{Appendix A}
If you need to include an appendix to support the research in your
thesis, you can place it at the end of the manuscript.

\backmatter
\listoffigures
\listoftables

\begin{acknowledgements}
  Here you might want to acknowledge someone.
\end{acknowledgements}

\end{document}
